% titel: Blatt 8, A2: Newton-Verfahren
% Teil II: der durchgerechnete Weg.

\subsection*{Gegeben und gesucht}
a) Ein Newton-Schritt für $f(x)=x^5-2x-1$ auf $[1,2]$ und Prüfung der
Bedingungen von Satz 4.3.8. b) Zwei Newton-Schritte für $e^{x}=2-x^2$ auf
$[0,1]$. c) Nullstelle $x^*\in[1,2]$ von $f$ aus (a) mit angepasstem
Python-Code (sehr genau).

\subsection*{a) Handschritt und Voraussetzungen}
\emph{Gedankenschritt:} $f'(x)=5x^4-2$ (Ableitung), Start $x_0=2$:
\begin{equation*}
  f(2)=32-4-1=27,\qquad f'(2)=5\cdot16-2=78,
\end{equation*}
\begin{equation*}
  x_1=2-\frac{27}{78}=2-0{,}3461538=1{,}6538462 .
\end{equation*}
\emph{Prüfung der Voraussetzungen (Satz 4.3.8):} $f(1)=-2<0<f(2)=27$;
$f'=5x^4-2>0$ auf $[1,2]$ (monoton steigend, Nullstelle eindeutig);
$f''=20x^3>0$ (konvex), Start $x_0=2$ auf der \emph{richtigen} Seite. Also
greifen die Bedingungen.

\subsection*{b) Zwei Newton-Schritte}
\emph{Gedankenschritt:} Umformen zu $g(x)=e^x+x^2-2=0$; der Newton-Schritt
braucht $g'(x)=e^x+2x$. Start $x_0=0{,}5$:
\begin{equation*}
  g(0{,}5)=e^{0{,}5}+0{,}25-2\approx1{,}6487-1{,}75=-0{,}1013,\qquad
  g'(0{,}5)=1{,}6487+1=2{,}6487,
\end{equation*}
\begin{equation*}
  x_1=0{,}5-\frac{-0{,}1013}{2{,}6487}=0{,}53824 .
\end{equation*}
Zweiter Schritt:
\begin{equation*}
  g(0{,}53824)\approx0{,}002683,\quad g'(0{,}53824)\approx2{,}7897,
  \quad x_2=0{,}53824-\frac{0{,}002683}{2{,}7897}=0{,}53728 .
\end{equation*}

\subsection*{c) Sehr genaue Nullstelle per Code}
\emph{Gedankenschritt:} Der Vorlesungscode wird nur leicht angepasst
(mehr Iterationen, Abbruch-bis-klein statt fester Schrittanzahl):
\begin{verbatim}
x = 2.0
while abs(x**5 - 2*x - 1) > 1e-30:
    x = x - (x**5 - 2*x - 1) / (5*x**4 - 2)
print(f"{x:.19f}")
\end{verbatim}
Ausgabe (Beispielwert entspricht $x^*$):
\begin{equation*}
  x^*\approx1{,}2906488013467097215\ldots
\end{equation*}

\begin{ergebnis}
a) $x_1=1{,}653846$; Bedingungen erfüllt. b) $x_1=0{,}53824$,
$x_2=0{,}53728$. c) $x^*\approx1{,}2906488013$ (sehr genau; 19 Nachkommastellen
gezeigt).
\end{ergebnis}

\begin{kontrolle}
\textbf{Probe (a):} $x_1=1{,}653846$ in $f$ ergibt $(1{,}653846)^5-2\cdot
1{,}653846-1\approx3{,}38$ -- deutlich kleiner als die anfängliche
Abweichung bei $x_0$; der Schritt ist die korrekte Tangentialkorrektur.
\checkmark

\textbf{Probe (b):} $e^{0{,}53728}+0{,}53728^2-2\approx0{,}0000$, also
Nullstelle (numerisch verifiziert); $x_2\approx0{,}53727$ ist konsistent
mit dem Residuum $\approx0$. \checkmark

\textbf{Grenzfall (quadratische Konvergenz):} ab Schritt 2 ändert sich
$x$ kaum noch; die Genauigkeit wächst quadratisch, weitere Schritte geben
fast keine zusätzlichen Dezimalen mehr. \checkmark
\end{kontrolle}

\begin{merksatz}
Newton: $x_{k+1}=x_k-\frac{f(x_k)}{f'(x_k)}$ mit $f'$ in der Nähe der
Nullstelle ungleich $0$. Vor dem Start Vorzeichenwechsel, Monotonie und
Krümmung prüfen. Im Code braucht man nur $f$ und $f'$ sowie eine
Abbruchbedingung (bis $|f(x_k)|$ klein).
\end{merksatz}
